% Irodov Part III §3.2, problems 66--77, PDF pp.109--110,
% printed pp.100--101. Fig.3.8 and Fig.3.9 checked against page images.
\electroconductproblem{66}{100}
图3.8中四块相邻大金属板
从上到下编号1、2、3、4，
相邻间距均为 $d$。
外侧板1、4用导线连通，
内侧2、3分别接电源正负端，
其电势差为 $\Delta\varphi>0$。
求 (a) 相邻空隙电场；
(b) 每块板的\emph{总面电荷密度}。
\exsolution
记 $\varphi_2-\varphi_3=\Delta\varphi$，
外侧共同电势
$\varphi_1=\varphi_4=\varphi_0$。
外侧板总电荷守恒且起初无净电荷，
故两外板净面密度之和为0：
\[
\frac{\varepsilon_0}{d}(\varphi_0-\varphi_2)
+\frac{\varepsilon_0}{d}(\varphi_0-\varphi_3)=0
\quad\Longrightarrow\quad
\varphi_0=\frac{\varphi_2+\varphi_3}{2}.
\]
以向下为正方向，
各间隙场分量为
$E_{i,i+1}=(\varphi_i-\varphi_{i+1})/d$：
\[
\boxed{E_{12}=-\frac{\Delta\varphi}{2d},\quad
E_{23}=\frac{\Delta\varphi}{d},\quad
E_{34}=-\frac{\Delta\varphi}{2d}}.
\]
各板总面电荷密度等于相邻表面电荷之和：
\[
\boxed{\sigma_1=-\frac{\varepsilon_0\Delta\varphi}{2d},
\quad \sigma_2=\frac{3\varepsilon_0\Delta\varphi}{2d},
\quad \sigma_3=-\frac{3\varepsilon_0\Delta\varphi}{2d},
\quad \sigma_4=\frac{\varepsilon_0\Delta\varphi}{2d}}.
\]
\excheck $\sigma_1+\sigma_4=0$、
$\sigma_2+\sigma_3=0$，
且中间间隙场强大小为两边的2倍。

\electroconductproblem{67}{100}
两块无限大平行导体板1、2相距 $\ell$。
一个点电荷 $q$ 处于板间，
距板1的距离为 $x$（$0<x<\ell$）。
求两板上\emph{总感应电荷}。
\exsolution
点电荷的场可用沿板平行平面的傅里叶分解、
镜像级数或 Green 函数处理。
计算\emph{总}感应电荷时，
不用求详细面密度：
设点电荷距板1为 $x$，
将点电荷想象沿平行板平面均匀铺展，
由于该等效化只改变表面电荷分布而不改变
两板\emph{总量}，
法向静电场在两段间距
$x$ 与 $\ell-x$ 上为分段常值。
导体板电势相同，
两段电场电势降必须相抵，
得到
\[
\boxed{Q_1=-q\,\frac{\ell-x}{\ell}},
\qquad
\boxed{Q_2=-q\,\frac{x}{\ell}}.
\]
\excheck $Q_1+Q_2=-q$，
且当 $x\to0$ 时板1感应电荷趋 $-q$；
原题英文正文有一处重复的距离语句，
以其最后明确的“距板1为 $x$”为准。

\electroconductproblem{68}{100}
任意形状导体表面局部面电荷密度为 $\sigma$。
求作用于导体单位表面积上的向外电压力。
\exsolution
在导体外表面，
总电场法向分量
$E_n=\sigma/\varepsilon_0$，
导体内部场为零。
但计算面电荷对自身的净力时，
须排除它的自场；
局部电荷感受到其余电荷产生的平均场
\[
E_{\rm other}=
\frac{E_{\rm outside}+E_{\rm inside}}2
=\frac{\sigma}{2\varepsilon_0}.
\]
所以单位面积受力模
\[
\boxed{P_e=\frac FS=
\sigma E_{\rm other}=
\frac{\sigma^2}{2\varepsilon_0}}.
\]
\excheck $P_e$ 是\emph{正的向外电压力}，
即使 $\sigma<0$，平方后仍为向外张力；
用总场乘 $\sigma$ 会错多出因子2。

\electroconductproblem{69}{101}
一颗半径 $R=1.5\,\mathrm{cm}$ 的带电金属球
总电荷 $q=10\,\mathrm{\mu C}$。
求作用在\emph{其中一个半球表面上的电荷}
所受合电力的模。
\exsolution
球面面电荷密度均匀：
$\sigma=q/(4\pi R^2)$。
第68题给单位面积向外电压力
$P_e=\sigma^2/(2\varepsilon_0)$。
在半球表面，
横向力由对称性抵消，
轴向分量积分
\[
F_z=P_e
\int_0^{2\pi}\!\!\int_0^{\pi/2}
\cos\theta R^2\sin\theta\,\mathrm d\theta\,\mathrm d\phi
=P_e\pi R^2.
\]
因此
\[
\boxed{F=\frac{\sigma^2\pi R^2}{2\varepsilon_0}
=\frac{q^2}{32\pi\varepsilon_0R^2}}.
\]
代入 SI 单位得
\[
\boxed{F\simeq5.0\times10^2\,\mathrm N=0.50\,\mathrm{kN}}.
\]
\excheck 不是简单的 $P_e$ 乘半球表面积 $2\pi R^2$；
因为每个面元受力方向不同，
应投影到共同的对称轴积分。

\electroconductproblem{70}{101}
一个原先中性的导体球，
半径为 $R$，
放在外加均匀电场中后，
其表面出现电荷密度
$\sigma(\theta)=\sigma_0\cos\theta$。
求所有同号诱导电荷所受的合静电力模。
\exsolution
选正半球
$0\le\theta\le\pi/2$，
每处向外电压力为
$P_e(\theta)=\sigma(\theta)^2/(2\varepsilon_0)
=\sigma_0^2\cos^2\theta/(2\varepsilon_0)$。
合力的横向分量抵消，
沿场轴积分：
\[
F=\frac{\sigma_0^2}{2\varepsilon_0}
2\pi R^2\int_0^{\pi/2}
\cos^3\theta\sin\theta\,\mathrm d\theta.
\]
积分为 $1/4$，
故
\[
\boxed{F=\frac{\pi R^2\sigma_0^2}{4\varepsilon_0}}.
\]
\excheck 因只有一个半球同号，
积分限是 $0$ 至 $\pi/2$；
力密度还需投影一次 $\cos\theta$。

\electroconductproblem{71}{101}
强度为
$E=1.0\,\mathrm{kV/cm}$ 的电场
使水分子发生极化。
若此极化等效于每 $N$ 个水分子中
有一个完全沿场方向取向，
求 $N$。
水分子电偶极矩
$p=0.62\times10^{-29}\,\mathrm{C\,m}$。
\exsolution
水的极化强度（线性介质近似）
$P=\varepsilon_0(\varepsilon_r-1)E$。
取室温水的介电常数
$\varepsilon_r\simeq81$、
水的数密度
$n_0=(\rho/M)N_A$，
$\rho\simeq10^3\,\mathrm{kg/m^3}$、
$M\simeq0.018\,\mathrm{kg/mol}$。
若每 $N$ 个分子仅一个完全取向，
等效极化
$P=n_0p/N$。
所以
\[
\boxed{N=\frac{n_0p}
{\varepsilon_0(\varepsilon_r-1)E}}.
\]
代入 $E=10^5\,\mathrm{V/m}$：
\[
\boxed{N\simeq3\times10^3}.
\]
\excheck 原题未明列水的介电常数、密度，
数值须采用室温物性近似；
极化“等效取向分数”不是水分子真正
完全平行排列的比例。

\electroconductproblem{72}{101}
一个不带永久偶极矩、但可极化率为 $\beta$
的非极性分子，距永久偶极矩大小为 $p$
的极性分子 $\ell$，
两个分子连线与 $\bm p$ 同轴。
求两分子间的作用力模。
\exsolution
按本版的\emph{体积极化率}记号，
诱导偶极矩
$\bm p_{\rm ind}=\varepsilon_0\beta\bm E$。
极性分子沿偶极轴的场强为
\[
E(\ell)=\frac{2p}{4\pi\varepsilon_0\ell^3}
=\frac{p}{2\pi\varepsilon_0\ell^3}.
\]
诱导偶极子的势能需包含极化建立功，
故
$U_{\rm ind}=-\frac12\varepsilon_0\beta E^2$。
力的模为
\[
F=\left|-\frac{\mathrm dU_{\rm ind}}{\mathrm d\ell}\right|
=\frac12\varepsilon_0\beta
\left|\frac{\mathrm d(E^2)}{\mathrm d\ell}\right|
=\frac{3\varepsilon_0\beta E^2}{\ell}.
\]
代入轴向电场：
\[
\boxed{F=\frac{3\beta p^2}
{4\pi^2\varepsilon_0\ell^7}}.
\]
\excheck 本题力为吸引。
若使用另一种极化率单位约定
$\bm p_{\rm ind}=\beta_{\rm SI}\bm E$，
则 $\beta_{\rm SI}=\varepsilon_0\beta$，
不能混淆两者。

\electroconductproblem{73}{101}
一个可极化的非极性分子
在半径 $R$、总电荷 $q$ 的均匀带电细圆环轴上，
距环心为 $x$。
求沿轴的作用力为零的位置、
两侧作用力模的极大位置，
并说明 $F_x(x)$ 的图像。
\exsolution
轴上环电场
$E_x=kqx/(R^2+x^2)^{3/2}$。
诱导偶极子的能量
$U=-\frac12\varepsilon_0\beta E_x^2$，
所以
\[
F_x=-\frac{\mathrm dU}{\mathrm dx}
=\frac12\varepsilon_0\beta
\frac{\mathrm d}{\mathrm dx}E_x^2
=\varepsilon_0\beta k^2q^2
\frac{x(R^2-2x^2)}{(R^2+x^2)^4}.
\]
因此当 $x\ge0$ 时，
\[
\boxed{F_x=0\ \Longleftrightarrow\
x=0\ \text{或}\ x=R/\sqrt2}.
\]
在非零位置 $x=R/\sqrt2$，
场强 $E$ 极大而力为零。
设 $u=x/R$，
力的无量纲形状
$f(u)=u(1-2u^2)/(1+u^2)^4$。
求导得到
$f'(u)\propto1-13u^2+10u^4$，
所以两个非零极值位置
\[
\boxed{\frac{x}{R}=
\sqrt{\frac{13-\sqrt{129}}{20}}
\simeq0.29\quad\text{（向外）}},
\]
\[
\boxed{\frac{x}{R}=
\sqrt{\frac{13+\sqrt{129}}{20}}
\simeq1.10\quad\text{（向内）}}.
\]
\excheck $x=0$ 因对称性也是零力点；
原书(a)只列非平凡零力点 $R/\sqrt2$。
力曲线为奇函数，$x>0$ 在
$0<x<R/\sqrt2$ 向外，
更远处向内，并最终按 $x^{-5}$ 衰减：
因圆环远轴场 $E_x\sim kq/x^2$，
$U_{\rm ind}=-\varepsilon_0\beta E_x^2/2\sim-x^{-4}$，
故
$F_x=-\mathrm dU_{\rm ind}/\mathrm dx\sim-x^{-5}$。
这里不能把\emph{点偶极场}的 $x^{-3}$
错当成带电圆环总电荷的远场。

\electroconductproblem{74}{101}
点电荷 $q$ 位于半径为 $R$、
相对介电常数为 $\varepsilon_r$
的均匀各向同性介质球心。
求介质内的极化强度 $\bm P(\bm r)$，
及以球心为中心、半径 $r<R$
的球内\emph{束缚电荷总量}。
\exsolution
由位移矢量高斯定律
$\oint\bm D\cdot\mathrm d\bm S=q$，
\[
\bm D(r)=\frac{q}{4\pi r^2}\bm e_r,\qquad
\bm E=\frac{\bm D}{\varepsilon_0\varepsilon_r}.
\]
极化强度
$\bm P=\varepsilon_0(\varepsilon_r-1)\bm E$，
所以
\[
\boxed{\bm P(\bm r)=
\frac{\varepsilon_r-1}{\varepsilon_r}
\frac{q\,\bm r}{4\pi r^3}}.
\]
束缚电荷净量
$Q'_b=-\oint\bm P\cdot\mathrm d\bm S$，
\[
\boxed{Q'_b=
-q\,\frac{\varepsilon_r-1}{\varepsilon_r}},
\qquad 0<r<R.
\]
\excheck 束缚电荷净量集中在点电荷附近的
分布奇性中；
在 $0<r<R$ 的其余均匀区域，
$\nabla\cdot\bm P=0$。

\electroconductproblem{75}{101}
证明导体与均匀各向同性电介质分界面上，
电介质束缚面电荷密度
$\sigma'_b=-\sigma(\varepsilon_r-1)/\varepsilon_r$，
其中 $\sigma$ 是导体自由电荷面密度。
\exsolution
静电平衡的导体内部 $\bm E=0$，
导体—介质法向位移边界条件给
\[
D_n=\sigma,
\quad D_n=\varepsilon_0\varepsilon_r E_n.
\]
所以介质法向电场
$E_n=\sigma/(\varepsilon_0\varepsilon_r)$。
介质极化为
$\bm P=\varepsilon_0(\varepsilon_r-1)\bm E$，
其在从导体指向介质法向的分量为
$P_n=\sigma(\varepsilon_r-1)/\varepsilon_r$。
但电介质的\emph{外法向}指向导体，
与上述方向相反，
故束缚面电荷
$\sigma'_b=\bm P\cdot\bm n_{\rm diel,out}=-P_n$，
得到
\[
\boxed{\sigma'_b=
-\sigma\frac{\varepsilon_r-1}{\varepsilon_r}}.
\]
\excheck 必须区分“导体外法向”
与“电介质外法向”；
它们在公共界面上方向相反。

\electroconductproblem{76}{101}
一个任意形状导体带净自由电荷 $q$，
其外面包围有一层均匀相对介电常数
$\varepsilon_r$ 的有限厚电介质，
如图3.9。
求介质的内、外两表面上的束缚电荷总量。
\exsolution
导体对周围介质的位移通量为 $q$。
均匀介质中
$\bm P=(\varepsilon_r-1)\bm D/\varepsilon_r$。
内边界按电介质外法向指向导体，
所以
\[
Q'_{\rm inner}=
-\frac{\varepsilon_r-1}{\varepsilon_r}
\oint_{\rm inner}\bm D\cdot\mathrm d\bm S_{\rm conductor}
=\boxed{-q\frac{\varepsilon_r-1}{\varepsilon_r}}.
\]
介质整体无体束缚电荷
（介电常数均匀且介质内部无外加自由体电荷），
$\nabla\cdot\bm P=0$。
所以外表面束缚电荷抵消内表面总量：
\[
\boxed{Q'_{\rm outer}=
+q\frac{\varepsilon_r-1}{\varepsilon_r}}.
\]
\excheck 两表面的总束缚电荷相加为零；
这是有限、均匀且无自由体电荷的介质层。

\electroconductproblem{77}{101}
相对介电常数为 $\varepsilon_r$ 的均匀介质
做成内半径 $a$、外半径 $b>a$ 的球壳。
球壳上存在一定正自由电荷，分别按两种方式：
(a) 均匀分布在内球面；
(b) 均匀分布在球壳体积。
定性给出电场 $E(r)$ 和电势 $\varphi(r)$ 的图线，
并写出分段表达式（总自由电荷记为 $Q$）。
\exsolution
设球壳外部及内腔均为真空，
取无穷远电势为零，
$k=1/(4\pi\varepsilon_0)$。
(a) 自由电荷全在 $r=a$ 球面。
球对称的高斯定律给
\[
\boxed{E_a(r)=
\begin{cases}
0,&r<a,\\
kQ/(\varepsilon_r r^2),&a<r<b,\\
kQ/r^2,&r>b.
\end{cases}}
\]
积分并要求电势连续：
\[
\boxed{\varphi_a(r)=
\begin{cases}
kQ\!\left[\dfrac1{\varepsilon_r a}
+\dfrac{1-1/\varepsilon_r}{b}\right],&r<a,\\
kQ\!\left[\dfrac1{\varepsilon_r r}
+\dfrac{1-1/\varepsilon_r}{b}\right],&a\le r\le b,\\
kQ/r,&r\ge b.
\end{cases}}
\]
(b) 体自由电荷密度
$\rho_f=3Q/[4\pi(b^3-a^3)]$。
球壳内部被包围自由电荷
$Q_{\rm enc}(r)=(4\pi\rho_f/3)(r^3-a^3)$，
故
\[
\boxed{E_b(r)=
\begin{cases}
0,&r<a,\\
\dfrac{\rho_f}{3\varepsilon_0\varepsilon_r}
\left(r-\dfrac{a^3}{r^2}\right),&a\le r\le b,\\
kQ/r^2,&r\ge b.
\end{cases}}
\]
并有
\[
\boxed{\varphi_b(r)=
\begin{cases}
\varphi_b(a),&r<a,\\
\dfrac{kQ}{b}+
\dfrac{\rho_f}{3\varepsilon_0\varepsilon_r}
\left[\dfrac{b^2-r^2}{2}
+a^3\left(\dfrac1b-\dfrac1r\right)\right],
&a\le r\le b,\\
kQ/r,&r\ge b.
\end{cases}}
\]
其中 $\varphi_b(a)$ 取中段 $r=a$ 的值。
图像上：
(a) $E$ 在 $a,b$ 两界面有跳跃，
介质段按 $1/r^2$ 衰减；
(b) $E$ 从 $r=a$ 的零连续上升，
至 $r=b$ 穿出介质时出现介电跳跃，
再按 $1/r^2$ 衰减；
两种情况下 $\varphi(r)$ 均连续，
且内腔为常势。
\excheck 介质边界若无自由面电荷，
连续的是 $D_n$，
不是 $E_n$；而在(a)的内球面
存在自由面电荷，$D_n$ 必有相应跃变。
